% ======================================================================
%  sec15.tex —— 第 15 节 应用：闭浸入的分类
%  底本：原书 p.77（印刷页 = PDF 页）
%  原书标题：Application: classification of closed immersions
% ======================================================================

\section{应用：闭浸入的分类}\label{sec:15}

\subsection{分类}\label{sec:15.1}

第 \ref{sec:14.2} 小节把相对谱 $\Spec_X$ 定义为拟凝聚 $\cO_X$-代数范畴上的函子。本节只用其中最特殊的一类代数：商代数 $\cO_X/\cI$。结论是，$X$ 的闭浸入与 $\cO_X$ 的拟凝聚理想是同一批对象的两种描述。

\begin{theorem}\label{thm:15.1.1}
设 $X$ 是概形。下述两个函子有定义，且互为逆函子：
\[
\begin{array}{rcl}
\{\text{到 $X$ 的闭浸入}\}^{\op} & \longleftrightarrow & \{\cO_X \text{ 的拟凝聚理想}\}\\[0.8ex]
(i \colon Y \to X) & \longmapsto & \ker(\cO_X \to i_*\cO_Y)\\[0.4ex]
\Spec_X(\cO_X/\cI) & \longmapsfrom & \cI \subseteq \cO_X.
\end{array}
\]
\end{theorem}

\begin{proof}
先证有定义。只需说明：对任意拟凝聚理想 $\cI$，结构态射 $\Spec_X(\cO_X/\cI) \to X$ 是闭浸入。由习题 \ref{exc:14.2.6}，这一断言在 $X$ 上局部，故可设 $X = \Spec(A)$ 是仿射概形。取理想 $I \subseteq A$，使 $\widetilde{I} \simeq \cI$。按构造有
\[
\Spec_X(\cO_X/\cI) \simeq \Spec(A/I),
\]
而结构态射 $\Spec(A/I) \to \Spec(A)$ 对应典范同态 $A \to A/I$，后者确是闭浸入。

再证两函子互逆。设 $\cI \subseteq \cO_X$ 是 $\cO_X$ 的拟凝聚理想，考虑 $i \colon Y = \Spec_X(\cO_X/\cI) \to X$。按定义，有拟凝聚 $\cO_X$-代数的典范同构 $\cO_X/\cI \simeq i_*\cO_Y$，于是
\[
\ker(\cO_X \to i_*\cO_Y) \simeq \ker(\cO_X \to \cO_X/\cI) \simeq \cI.
\]
也就是说，先取左向函子、再取右向函子，所得与恒等函子自然同构。

另一方面，设 $i \colon Y \to X$ 是闭浸入。按 $\Spec_X$ 的定义，恒等同态 $i_*\cO_Y \to i_*\cO_Y$ 对应一个 $X$-态射
\[
\phi \colon Y \to \Spec_X(i_*\cO_Y).
\]
只需说明 $\phi$ 是同构。按构造，$\phi$ 在相应的定义理想之间诱导同构；于是由引理 \ref{lem:11.4.5}，$\phi$ 即为所求的同构。
\end{proof}

\begin{yizhu}
定理中的「拟凝聚」四字不能省。$\Spec_X$ 只在拟凝聚 $\cO_X$-代数上有定义（原书推论 \ref{cor:14.2.5}），若 $\cI$ 只是一般的理想层，$\cO_X/\cI$ 未必拟凝聚，$\Spec_X(\cO_X/\cI)$ 根本无从谈起。换言之，定理右端的「理想」必须限定在拟凝聚者之内，否则函子连定义域都不完整。警告 \ref{warn:16.2.4} 给的正是这种理想层。
[依据：原书 p.75（$\Spec_S$ 只接受拟凝聚 $\cO_S$-代数）与 p.80 警告 \ref{warn:16.2.4}。]
\end{yizhu}

\begin{yizhu}
左侧取相反范畴，是因为函子在态射上反向。闭浸入之间的 $X$-态射 $f \colon Y \to Y'$ 给出 $\cO_X$-代数的满射 $\cO_X/\cI_{Y'} \to \cO_X/\cI_Y$，即理想的反向包含 $\cI_Y \subseteq \cI_{Y'}$；而 $\Spec_X$ 本身又把 $\cO_X$-代数送到 $(\mathsf{Sch}'_X)^{\op}$（推论 \ref{cor:14.2.5}）。两次取反相互抵消，定理才成为范畴层面的等价，而非仅仅是对象之间的双射。
[依据：原书 p.77 定理 \ref{thm:15.1.1} 与 p.75 推论 \ref{cor:14.2.5}。]
\end{yizhu}

\begin{corollary}\label{cor:15.1.2}
设 $X = \Spec(A)$ 是仿射概形。到 $X$ 的任意闭浸入都同构于某个 $\Spec(A/I) \to \Spec(A)$，其中 $I \subseteq A$ 是唯一确定的理想。
\end{corollary}

\begin{yizhu}
仿射情形可以把定理翻译成一句交换代数：$\Spec(A)$ 的闭子概形与 $A$ 的理想一一对应。同一定理在 Stacks Project 中分成两半陈述：闭浸入的定义理想是拟凝聚理想，且在仿射开集 $U = \Spec(R)$ 上该理想限制为 $\widetilde{I}$，$i^{-1}(U) \to U$ 就是 $\Spec(R/I) \to \Spec(R)$。
[来源：Stacks Project, Tag 01IN（Lemma 26.10.1，闭浸入的局部形状与定义理想），\url{https://stacks.math.columbia.edu/tag/01IN}（2026-11-17 访问）。]
\end{yizhu}

\begin{corollary}\label{cor:15.1.3}
闭浸入对基变换稳定。
\end{corollary}

\begin{proof}
设 $i \colon Y \to X$ 是 $S$-概形之间的闭浸入，其定义理想为 $\cI_Y \subseteq \cO_X$。把 $Y$ 与 $\Spec_X(\cO_X/\cI_Y)$ 等同。对任意 $S' \to S$，置 $X' := X \times_S S'$ 与 $Y' := Y \times_S S'$。则
\[
Y' \simeq \Spec_X(\cO_X/\cI_Y) \times_S S' \simeq \Spec_X(\cO_X/\cI_Y) \times_X X' \simeq \Spec_{X'}\bigl(f^*(\cO_X/\cI_Y)\bigr),
\]
其中 $f \colon X' \to X$ 是典范投影。注意 $f^*(\cO_X/\cI_Y) \simeq \cO_{X'}/f^*\cI_Y$，而 $f^*\cI_Y$ 是 $\cO_{X'}$ 的拟凝聚理想。故 $Y' \to X'$ 是闭浸入。
\end{proof}

\begin{yizhu}
证明中三处同构各有出处，值得逐一核对。第一处是 $Y \simeq \Spec_X(\cO_X/\cI_Y)$，即定理 \ref{thm:15.1.1}；第二处把 $\times_S S'$ 换成 $\times_X X'$，因为 $X' = X \times_S S'$ 配着同一投影；第三处是 $\Spec$ 与基变换的交换律，即习题 \ref{exc:14.2.6}，它要求 $f^*\cI_Y$ 拟凝聚，而拟凝聚层对拉回封闭。末步用 $f^*$ 右正合，把商层拉回成商层。
[依据：原书 p.77 推论 \ref{cor:15.1.3} 的证明与 p.75 习题 \ref{exc:14.2.6}。]
\end{yizhu}

结合习题 \ref{exc:10.1.4}，可得

\begin{corollary}\label{cor:15.1.4}
局部闭浸入对基变换稳定。
\end{corollary}

\begin{yizhu}
局部闭浸入按定义是开浸入与闭浸入的复合（习题 \ref{exc:10.1.4}）。开浸入对基变换的稳定性是纤维积的泛性质：$U \times_X X' \to X'$ 仍是开浸入。把两种稳定性合起来，复合的稳定性便随之得到，无需再对局部闭的情形单独验证。
[依据：原书 p.77 推论 \ref{cor:15.1.4} 与 p.59 习题 \ref{exc:10.1.4}。]
\end{yizhu}
